\chapter{Analytic construction of the functionals of the triadic tower}
\label{chap:funcionales-espectrales-triadicos}

The preceding synthesis has placed the spectral constants within the
triadic tower. This block contains their primary analytic construction:
a sequence of self-adjoint operators on cycles of order \(3^m\), a limit of
heat traces and a Mellin transform. The classical constants appear
as evaluations, finite parts or derivatives of the resulting meromorphic
functional, not as prefixes recognized in a catalogue.

The normalization of the operator $D_m$ explicitly contains $\pi$ in the scale
relating the discrete Laplacian to the Gaussian heat kernel. The
subsection therefore constitutes an analytic extension of the Archimedean
branch already constructed. The subsequent spectral constants are obtained
simultaneously from the resulting functional, with that normalization as a starting
hypothesis.

\section{Inductive limits and refinement scale}

The spectral extension requires distinguishing two analytic levels. The
first corresponds to bounded compatible families; the second to
unbounded operators, for which algebraic intertwining does not
by itself resolve domain problems.

\begin{theorem}[Bounded inductive limit]
\label{pf84:E03-06}
Let \(J_m:\mathcal H_m\to\mathcal H_{m+1}\) be isometries and
\(A_m\in\mathcal B(\mathcal H_m)\) self-adjoint operators such that
\[
 A_{m+1}J_m=J_mA_m,
 \qquad
 \sup_m\lVert A_m\rVert<\infty.
\]
There exists a unique bounded self-adjoint operator \(A\) on
\(\mathcal H_\infty=\varinjlim(\mathcal H_m,J_m)\) whose restriction to each
level is \(A_m\).
\end{theorem}

\begin{proof}
On the dense algebraic union we define
\[
 A[J_{m,\infty}\xi]:=J_{m,\infty}A_m\xi.
\]
Intertwining makes the definition independent of the representative. The
uniform bound allows \(A\) to be extended by continuity to
\(\mathcal H_\infty\); symmetry on the dense subspace and boundedness
imply its self-adjointness.
\end{proof}

\begin{remark}[Control for unbounded operators]
\label{pf84:E03-07}
If the \(A_m\) are unbounded, the identity
\(A_{m+1}J_m=J_mA_m\) only defines an action on a union of domains.
To obtain a closed or self-adjoint operator, the transported domains must be
declared, closability proved and deficiency indices
controlled, or essential self-adjointness proved on a compatible common
core. Omitting these obligations constitutes a type error, not an
abbreviation of the proof.
\end{remark}

Now let \(\mathcal S_\infty\) be a nonempty set of compatible surviving
sections and let \(\mathcal P_m\) be the finite set of their
prefixes of length \(m\). We write
\[
 \mathcal H_m:=\ell^2(\mathcal P_m),
 \qquad
 E_m(\gamma,\gamma')=
 \begin{cases}
 1,&\gamma'\text{ extends }\gamma,\\
 0,&\text{otherwise},
 \end{cases}
\]
and
\[
 d_m(\gamma):=
 \#\{\gamma'\in\mathcal P_{m+1}:\gamma'\mapsto\gamma\},
 \qquad
 D_m^{\deg}:=E_mE_m^*=\operatorname{diag}(d_m(\gamma)).
\]
We restrict ourselves to the live subspace \(d_m(\gamma)>0\), and each prefix of
level \(m+1\) has a unique parent.

\begin{theorem}[Normalized extension isometry]
\label{pf84:C01}
The operator
\[
 J_m:=E_m^*(D_m^{\deg})^{-1/2}:
 \mathcal H_m\longrightarrow\mathcal H_{m+1}
\]
is an isometry, and
\[
 J_me_\gamma=
 \frac1{\sqrt{d_m(\gamma)}}
 \sum_{\gamma'\mapsto\gamma}e_{\gamma'}.
\]
\end{theorem}

\begin{proof}
Uniqueness of the parent gives
\[
 J_m^*J_m
 =(D_m^{\deg})^{-1/2}E_mE_m^*(D_m^{\deg})^{-1/2}=I.
\]
The formula on \(e_\gamma\) is obtained by expanding the corresponding
column of \(E_m^*\).
\end{proof}

Setting
\[
 \mathcal K_m:=\mathcal H_{m+1}\ominus J_m\mathcal H_m,
 \qquad
 \mathcal H_{\mathrm{surv}}
 \cong\mathcal H_0\oplus\bigoplus_{m\ge0}\mathcal K_m,
\]
the refinement operator is defined
\[
 \mathcal R|_{\mathcal H_0}=0,
 \qquad
 \mathcal R|_{\mathcal K_m}=3^{m+1}I_{\mathcal K_m},
\]
with maximal domain
\[
 \Dom(\mathcal R)=
 \mathcal H_0\oplus
 \left\{(\xi_m)_{m\ge0}:
 \sum_{m\ge0}3^{2m+2}\lVert\xi_m\rVert^2<\infty\right\}.
\]

\begin{theorem}[Refinement operator]
\label{pf84:C02}
If each \(\mathcal K_m\) is finite-dimensional, \(\mathcal R\) is positive,
self-adjoint and has compact resolvent. Its spectrum measures refinement
depth; it is not a spectrum of primes or zeros.
\end{theorem}

\begin{proof}
The orthogonal decomposition reduces \(\mathcal R\) to real scalar blocks
on its maximal domain, which proves positivity and self-adjointness. The
eigenvalues \(3^{m+1}\) escape to infinity and their multiplicities are finite;
therefore the resolvent is compact. An infinite-dimensional block with
bounded eigenvalue would falsify this last conclusion.
\end{proof}

\section{Integer, prime and orbital operators}

Let \(\mathscr W_{\mathrm{bal}}\) be the set of balanced ternary
words reduced by theorem
\cref{pf84:C03}. On \(\ell^2(\mathscr W_{\mathrm{bal}})\), the operator
\[
 Be_w=b(w)e_w
\]
is taken on its maximal domain. Its restriction to the subspace of positive
integers is denoted by \(Q\).

\begin{theorem}[Integer and positive operators]
\label{pf84:C04}
The operator \(B\) is self-adjoint, has simple spectrum \(\ZZ\) and
compact resolvent. The operator \(Q\) has simple spectrum
\(\NN_{\ge1}\) and compact resolvent.
\end{theorem}

\begin{proof}
The bijection of \cref{pf84:C03} enumerates each integer exactly once. A
real diagonal operator on its maximal domain is self-adjoint; since its entries
escape to infinity outside finite sets, its resolvent is compact. The
positive restriction preserves these properties.
\end{proof}

The omitted classical proof of theta inversion, continuation and the functional
equation is merged only once with the triadic realization already constructed
in \cref{pf84:C05}. There,
\(\Theta_B=1+2\Theta_+\) and \(Z=\mathcal Z_3\) are fixed literally; therefore this block does not
define a second origin of \(\zeta\).

Let \(\mathbb P\) be the set of ordinary primes,
\(\mathfrak h_{\mathbb P}:=\ell^2(\mathbb P)\) and
\[
 H_{\mathbb P}e_p=(\log p)e_p.
\]
In the bosonic Fock space \(\Gamma_s(\mathfrak h_{\mathbb P})\), the occupations
\((k_p)_p\) have finite support.

\begin{theorem}[Fock representation of factorization]
\label{pf84:C08}
The map on bases
\[
 U_F\lvert(k_p)_p\rangle=e_{\prod_pp^{k_p}}
\]
extends to a unitary operator
\[
 U_F:\Gamma_s(\mathfrak h_{\mathbb P})
 \longrightarrow\ell^2(\NN_{\ge1})
\]
and satisfies
\[
 U_F\,d\Gamma(H_{\mathbb P})\,U_F^*=\log Q.
\]
\end{theorem}

\begin{proof}
Unique factorization gives a bijection between finite occupations and positive
integers, hence a unitary identification of the orthonormal bases.
Moreover,
\[
 \sum_pk_p\log p=\log\!\left(\prod_pp^{k_p}\right),
\]
which proves the intertwining. The explicit starting structure is the
choice of prime modes and symmetric Fock space.
\end{proof}

On
\[
 \mathcal H_{\mathrm{orb}}
 :=\ell^2\{(p,k,\varepsilon):
 p\in\mathbb P,\ k\ge1,\ \varepsilon=\pm1\},
\]
we define
\[
 D_{\mathrm{orb}}e_{p,k,\varepsilon}
 =\varepsilon k\log p\,e_{p,k,\varepsilon},
 \qquad
 W_{\log}e_{p,k,\varepsilon}
 =(\log p)e_{p,k,\varepsilon}.
\]

\begin{proposition}[Orbital trace]
\label{pf84:C09}
The operator \(D_{\mathrm{orb}}\) is self-adjoint and has compact resolvent. If
\(f\in C_c(\RR)\) is bounded, then
\[
 \Tr\!\left(
 W_{\log}e^{-|D_{\mathrm{orb}}|/2}f(D_{\mathrm{orb}})
 \right)
 =
 \sum_{p,k}(\log p)p^{-k/2}
 \bigl[f(k\log p)+f(-k\log p)\bigr].
\]
\end{proposition}

\begin{proof}
The diagonal is real and its lengths escape every compact set with
finite multiplicity. The compact support of \(f\) makes the regulated
product finite-rank; its trace is the sum of the diagonal entries.
The index \(p\) here ranges over ordinary primes. The result realizes the
classical prime side, but does not identify \(D_{\mathrm{orb}}\) with an operator
dual to the zeros.
\end{proof}

\section{Completed Guinand--Weil form}

We fix
\[
 \mathcal D_{\mathrm W}:=C_c^\infty(\RR),
 \qquad
 \widehat\psi(t)=\int_{\RR}\psi(x)e^{-itx}\,\dd x,
 \qquad
 \psi(x)=\frac1{2\pi}\int_{\RR}\widehat\psi(t)e^{itx}\,\dd t,
\]
and, for \(\psi,\phi\in\mathcal D_{\mathrm W}\),
\[
 h_{\psi,\phi}(z)
 :=\widehat\psi(z)\,
 \overline{\widehat\phi(\overline z)}.
\]
On the real line \(h_{\psi,\psi}=|\widehat\psi|^2\).

\begin{definition}[Test class]
\label{pf84:E01}
The test class is \(\mathcal D_{\mathrm W}\), stable under differentiation,
reflection and translation. The inverse used is
\[
 \check h(x)=\frac1{2\pi}\int_{\RR}h(t)e^{itx}\,\dd t.
\]
Smoothness and compact support provide rapid decay in strips
and compact support for \(\check h\), conditions ensuring the
convergence of the following blocks.
\end{definition}

Sea \(\psi_\Gamma=\Gamma'/\Gamma\). Definimos
\[
 \Irr_\#:=
 \left\{\operatorname{ev}_9(w):
 w\in\mathscr W_9^+,\ 
 P_{\Irr,\#}e_w=e_w\right\}.
\]
By \cref{pf84:C07}, this image coincides with the ordinary primes; here it is
used only to type the index of the orbital block, without repeating that
proof.
\begin{align}
 \mathcal P(h)&:=h(i/2)+h(-i/2),\\
 \mathcal G(h)&:=-\log\pi\,\check h(0)
 +\frac1{2\pi}\int_{\RR}h(t)
 \Re\psi_\Gamma\!\left(\frac14+\frac{it}{2}\right)\dd t,\\
 \mathcal O_\#(\check h)&:=
 \sum_{p\in\Irr_\#}\sum_{k\ge1}
 (\log p)p^{-k/2}
 \bigl[\check h(k\log p)+\check h(-k\log p)\bigr].
\end{align}

\begin{definition}[Completed functional]
\label{pf84:E02}
For \(h=h_{\psi,\phi}\), we fix
\[
 \mathcal W(h):=\mathcal P(h)+\mathcal G(h)
 -\mathcal O_\#(\check h).
\]
The orbital block comes from the irreducible selector of the
\cref{chap:producto-nativo}; the polar and gamma blocks belong to the
classical completion of zeta. This difference in provenance is preserved in
the notation and the proof.
\end{definition}

\begin{theorem}[Guinand--Weil explicit formula]
\label{pf84:E03}
For \(\psi,\phi\in C_c^\infty(\RR)\) and \(h=h_{\psi,\phi}\),
\[
 \lim_{T\to\infty}
 \sum_{\substack{\rho:\,\zeta(\rho)=0\\|\Im\rho|\le T}}
 h\!\left(\frac{\rho-1/2}{i}\right)
 =\mathcal W(h),
\]
where the nontrivial zeros are counted with multiplicity and the limit is
taken with classical symmetrization.
\end{theorem}

\begin{proof}
Paley--Wiener places \(h_{\psi,\phi}\) in the admissible class of the explicit
formula. The poles of
\(\Lambda_\zeta(s)=\pi^{-s/2}\Gamma(s/2)\zeta(s)\) produce
\(\mathcal P\), the Archimedean factor produces \(\mathcal G\) and the Euler
product produces the sum over prime powers. The identification of this
last sum with \(\mathcal O_\#\) is made only through
\cref{pf84:E04}; thus the proof of native primality is not duplicated.
\end{proof}

\begin{proposition}[Reindexing of the orbital block]
\label{pf84:E04}
If the irreducibles of the native product coincide exactly with the ordinary
primes and the weights
\((\log p)p^{-k/2}\) and the two orientations \(\pm k\log p\) are preserved, then
\[
 \mathcal O_\#(\check h)=\mathcal O(\check h).
\]
\end{proposition}

\begin{proof}
The bijection between native irreducibles and primes reindexes the sum without
modifying indices, weights or arguments. A nonprime irreducible, a missing
prime or a change in weight would falsify the identity.
\end{proof}

\begin{remark}[Meromorphic--entire separation]
\label{pf84:E05}
The function \(\Lambda_\zeta\) has simple poles at \(0\) and \(1\), which
give rise to the polar block. In contrast,
\[
 \xi(s)=\frac12s(s-1)\Lambda_\zeta(s)
\]
is entire: the factors \(s(s-1)\) cancel precisely those poles. Therefore no
poles at \(0\) or \(1\) are attributed to \(\xi\).
\end{remark}

\begin{definition}[Hermitian form and preoperator]
\label{pf84:E06}
We define
\[
 \mathcal B_{\mathrm W}(\psi,\phi)
 :=\mathcal W(h_{\psi,\phi}),
 \qquad
 \mathcal V_{\mathrm W}:=
 \mathcal D_{\mathrm W}/\operatorname{Rad}\mathcal B_{\mathrm W},
\]
and, on the algebraic quotient, the candidate rule
\[
 D_{\mathrm W,0}[\psi]:=[-i\psi'].
\]
It is not assumed that \(\mathcal B_{\mathrm W}\) is positive or that
\(\mathcal V_{\mathrm W}\) is already a Hilbert space. The rule becomes
well-defined only after invariance of the radical has been proved.
\end{definition}

\begin{theorem}[Algebraic symmetry and invariances]
\label{pf84:E07}
The radical is invariant under \(-i\,\dd/\dd x\);
\(D_{\mathrm W,0}\) is well-defined and
\(\mathcal B_{\mathrm W}\)-symmetric. The translations
\(T_a\psi(x)=\psi(x-a)\) preserve the form and the reflection
\(J\psi(x)=\psi(-x)\) satisfies
\[
 JD_{\mathrm W,0}J=-D_{\mathrm W,0}.
\]
\end{theorem}

\begin{proof}
The Fourier convention gives
\[
 h_{-i\psi',\phi}(z)=h_{\psi,-i\phi'}(z).
\]
The identity holds before applying any of the three blocks of
\(\mathcal W\), so the derivative preserves the radical and the
preoperator is symmetric. The opposite translation phases cancel
in \(h\); reflection changes the sign of the derivative and leaves the
polar, gamma and orbital blocks invariant. This algebraic symmetry is not yet equivalent
to self-adjointness in a positive Hilbert space.
\end{proof}

\section{Local geometry of the critical line}

\begin{lemma}[Möbius coordinate of the critical line]
\label{pf84:RH-01}
For \(s\in\CC\setminus\{0,1\}\), let
\[
 \tau(s):=i\frac{s-1}{s}.
\]
With \(F(s)=1-s\) and \(C(s)=\overline s\),
\[
 \tau(Fs)=-\frac1{\tau(s)},
 \qquad
 \tau(Cs)=-\overline{\tau(s)},
\]
and
\[
 \tau(Fs)=\tau(Cs)
 \iff |\tau(s)|=1
 \iff \Re s=\frac12.
\]
\end{lemma}

\begin{proof}
The two identities are obtained by direct substitution. Their equality
is equivalent to \(\tau(s)^{-1}=\overline{\tau(s)}\), that is,
\(|\tau(s)|=1\); expanding this last condition gives
\(\Re s=1/2\). For the zeros of zeta, functional and conjugate symmetries
are also required. The equivalence parametrizes the line; it does not by
itself place the zeros on it.
\end{proof}

\begin{lemma}[Local Hermite--Biehler sign]
\label{pf84:RH-04}
Let \(A\) be a real entire function with a simple real zero \(a\). For
\(z=a+iy\), \(y>0\) sufficiently small,
\[
|A(z)-iA'(z)|^2-|A(z)+iA'(z)|^2
 =-4A'(a)^2y+O(y^2)<0.
\]
The local convention associated with \(A\) is
\[
 E_A:=A+iA'.
\]
Only by specializing \(A=\Xi\) and additionally assuming that \(a\) is a simple real
zero of \(\Xi\) do we obtain
\[
 E_\Xi=\Xi+i\Xi'.
\]
\end{lemma}

\begin{proof}
We expand \(A\) and \(A'\) in Taylor series around \(a\), using
\(A(a)=0\) and \(A'(a)\ne0\), and retain the first-order terms in
\(y\). The result is strictly local: it does not prove that \(E_A\) has the
global Hermite--Biehler property.
\end{proof}

\begin{proposition}[Obstruction by counting laws]
\label{pf84:E14}
Let
\[
 N_A(T):=\Tr\mathbf1_{[-T,T]}(A).
\]
The operators \(Q\), \(B\), \(D_{\mathrm{orb}}\) and \(\mathcal R\) do not,
by themselves, have the counting law required of a possible operator dual
to the zeros.
\end{proposition}

\begin{proof}
Para \(T\ge1\),
\[
 N_Q(T)=\lfloor T\rfloor,
 \qquad
 N_B(T)=2\lfloor T\rfloor+1.
\]
The prime number theorem gives
\[
 N_{D_{\mathrm{orb}}}(T)
 =2\sum_{p\le e^T}\left\lfloor\frac{T}{\log p}\right\rfloor
 \ge2\pi(e^T)\sim\frac{2e^T}{T}.
\]
Finally, \(N_{\mathcal R}(T)\) is constant on each multiplicative
interval \([3^m,3^{m+1})\), since \(\mathcal R\) only has levels
\(3^{m+1}\) with finite multiplicities.

For
\(\Xi(z)=\xi(1/2+iz)\), the functional equation implies
\(\Xi(-z)=\Xi(z)\); the zeros \(\pm\gamma\) have the same multiplicity and
\[
 N_{\mathrm{sym}}(T)=2N_+(T)+O(1).
\]
Consequently, the bilateral convention requires
\[
 N_{\mathrm{sym}}(T)\sim\frac{T}{\pi}\log T,
\]
whereas the equivalent unilateral one, under this symmetry and pairing of
multiplicities, requires
\[
 N_+(T)\sim\frac{T}{2\pi}\log T.
\]
None of the four preceding functions has that asymptotic behavior.
\end{proof}

\section{Laplacians on cycles of order \texorpdfstring{$3^m$}{powers of three}}

For $m\geq 1$, let

\[
 N_m=3^m,
 \qquad
 C_m=\ZZ/N_m\ZZ,
\]

and consider the combinatorial Laplacian

\[
 (\Delta_m f)(j)=2f(j)-f(j+1)-f(j-1),
 \qquad f\in\ell^2(C_m).
\]

The characters

\[
 e_{m,n}(j)=N_m^{-1/2}
 \exp\!\left(\frac{2\pi i n j}{N_m}\right),
 \qquad n\in C_m,
\]

form an orthonormal basis and satisfy

\[
 \Delta_m e_{m,n}
 =4\sin^2\!\left(\frac{\pi n}{N_m}\right)e_{m,n}.
\]

Since $N_m$ is odd, each class of $C_m$ has a unique representative
$\widetilde n$ with

\[
 -\frac{N_m-1}{2}\leq\widetilde n\leq\frac{N_m-1}{2}.
\]

We define the self-adjoint operator $D_m$ by

\begin{equation}
 D_m e_{m,n}
 =\frac{N_m}{\sqrt\pi}
 \sin\!\left(\frac{\pi\widetilde n}{N_m}\right)e_{m,n}.
 \label{eq:operador-espectral-Dm}
\end{equation}

Then

\begin{equation}
 D_m^2=\frac{N_m^2}{4\pi}\,\Delta_m.
 \label{eq:Dm-laplaciano}
\end{equation}

Let $P_0$ be the orthogonal projection onto the constant character. For every
bounded even function $F$ on the spectrum of $D_m$, we introduce the positive
trace

\begin{equation}
 \operatorname{Tr}^{+}_m\!\bigl(F(D_m)\bigr)
 :=\frac12\operatorname{Tr}
 \!\left(F(D_m)P_0^{\perp}\right).
 \label{eq:traza-positiva}
\end{equation}

The factor $1/2$ selects one contribution from each spectral pair
$\{n,-n\}$; removing $P_0$ excludes the zero eigenvalue.

\begin{theorem}[Theta limit of the heat traces]
\label{thm:limite-theta-triadico}
For every $x>0$,
\[
 \lim_{m\to\infty}
 \operatorname{Tr}^{+}_m\!\left(e^{-xD_m^2}\right)
 =\Theta_+(x)
 :=\sum_{n\geq1}e^{-\pi n^2x}.
\]
\end{theorem}

\begin{proof}
The bilateral trace, once the constant mode is removed, is
\[
 \sum_{\substack{n\in C_m\\n\neq0}}
 \exp\!\left[
 -\frac{N_m^2x}{\pi}
 \sin^2\!\left(\frac{\pi\widetilde n}{N_m}\right)
 \right].
\]
For each fixed integer $n$, the summand converges to $e^{-\pi n^2x}$. Moreover,
for $|\widetilde n|\leq N_m/2$, the inequality
$\sin y\geq 2y/\pi$, valid on $0\leq y\leq\pi/2$, provides a
summable Gaussian majorant independent of $m$. Dominated
convergence applied to the discrete sum implies
\[
 \lim_{m\to\infty}
 \operatorname{Tr}\!\left(e^{-xD_m^2}P_0^\perp\right)
 =\sum_{n\in\ZZ\setminus\{0\}}e^{-\pi n^2x}.
\]
The division by two established in \eqref{eq:traza-positiva} preserves one
contribution from each pair $\{n,-n\}$ and gives $\Theta_+(x)$.
\end{proof}

\section{Theta--Mellin transform}

For $\Re s>1$, define

\begin{equation}
 \mathcal Z_{3}(s)
 :=
 \frac{\displaystyle
 \int_0^\infty\Theta_+(x)x^{s/2-1}\,\dd x}
 {\pi^{-s/2}\Gamma(s/2)}.
 \label{eq:funcional-theta-mellin}
\end{equation}

Absolute convergence allows term-by-term integration. Since

\[
 \int_0^\infty e^{-\pi n^2x}x^{s/2-1}\,\dd x
 =\Gamma(s/2)(\pi n^2)^{-s/2},
\]

we obtain

\begin{equation}
 \mathcal Z_3(s)=\sum_{n\geq1}n^{-s}=\zeta(s),
 \qquad \Re s>1.
 \label{eq:Z3-zeta}
\end{equation}

The equality extends meromorphically through continuation of the Riemann
zeta function \cite{Apostol1976}. Henceforth, all evaluations
come from \eqref{eq:funcional-theta-mellin}, its derivatives or the
residual decomposition of its Dirichlet series.

\begin{theorem}[Theta inversion and functional equation]
\label{pf84:C05}
With the preceding operator notation,
\[
 \Theta_B(x):=\Tr(e^{-\pi xB^2})
 =1+2\Theta_+(x),
 \qquad
 Z(s):=\Tr(Q^{-s})=\mathcal Z_3(s)
 \quad(\Re s>1).
\]
Poisson summation and the Mellin transform produce
\[
 \Theta_B(x)=x^{-1/2}\Theta_B(1/x)
\]
and the continuation of
\[
 \xi(s)=\frac12s(s-1)\pi^{-s/2}\Gamma(s/2)Z(s),
 \qquad
 \xi(s)=\xi(1-s).
\]
\end{theorem}

\begin{proof}
The equality \(\Theta_B=1+2\Theta_+\) separates the zero mode and pairs
\(\pm n\); the identity \(Z=\mathcal Z_3\) is
\eqref{eq:Z3-zeta}. Poisson summation applied to the Gaussian gives the
inversion of \(\Theta_B\). In the half-plane of convergence, term-by-term
integration provides
\[
 \pi^{-s/2}\Gamma(s/2)Z(s)
 =\int_0^\infty x^{s/2-1}\Theta_+(x)\,\dd x.
\]
Splitting the integral into \((0,1)\) and \((1,\infty)\), applying theta
inversion to the first part and separating the two classical poles yields the
continuation and the functional equation. This single incorporated classical proof
does not locate the zeros of \(\zeta\).
\end{proof}

\section{Residual decomposition modulo three}

For $r\in\{0,1,2\}$, let

\begin{equation}
 Z_r(s)=
 \sum_{\substack{n\geq1\\n\equiv r\pmod3}}n^{-s}.
 \label{eq:canales-residuales-Zr}
\end{equation}

In terms of the Hurwitz zeta function,

\[
 Z_0(s)=3^{-s}\zeta(s),
 \qquad
 Z_1(s)=3^{-s}\zeta\!\left(s,\frac13\right),
 \qquad
 Z_2(s)=3^{-s}\zeta\!\left(s,\frac23\right).
\]

The three terms have residue $1/3$ at $s=1$. We write their
Laurent expansions in the form

\begin{equation}
 Z_r(1+z)-\frac1{3z}=\sum_{n\geq0}c_{r,n}\,z^n,
 \qquad
 \zeta(1+z)-\frac1z=\sum_{n\geq0}a_n\,z^n.
 \label{eq:coeficientes-Laurent-residuales}
\end{equation}

\begin{theorem}[Decomposition of the Laurent coefficients]
\label{thm:descomposicion-laurent-triadica}
For every $n\geq0$,
\[
 a_n=c_{0,n}+c_{1,n}+c_{2,n},
 \qquad
 \gamma_n=(-1)^n n!\,a_n.
\]
If $\chi_{-3}$ denotes the nonprincipal real character modulo three, then
\[
 \frac{L^{(n)}(1,\chi_{-3})}{n!}=c_{1,n}-c_{2,n}.
\]
Furthermore, with \(\lambda=\log3\),
\[
 c_{0,n}=\frac13\left[
 \frac{(-\lambda)^{n+1}}{(n+1)!}
 +\sum_{k=0}^{n}a_k\frac{(-\lambda)^{n-k}}{(n-k)!}
 \right].
\]
These three identities uniquely determine
\((c_{0,n},c_{1,n},c_{2,n})\) from
\((a_n,L^{(n)}(1,\chi_{-3}))\) and the aggregated coefficients of lower
order.
\end{theorem}

\begin{proof}
The partition of the positive integers into residue classes provides
$\zeta=Z_0+Z_1+Z_2$. Subtracting the poles at $s=1$ and comparing the
coefficients of \eqref{eq:coeficientes-Laurent-residuales} yields the
first identity. The second coincides with the usual convention for the
Stieltjes constants. Moreover,
$Z_1-Z_2=L(s,\chi_{-3})$; the poles of $Z_1$ and $Z_2$ cancel and
comparison of coefficients gives the third relation. Finally,
\[
 Z_0(1+z)=\frac13e^{-\lambda z}\zeta(1+z).
\]
The product of the two Laurent series provides the explicit formula
for \(c_{0,n}\). The sum and difference then recover
\(c_{1,n}\) and \(c_{2,n}\), proving the asserted uniqueness.
\end{proof}

The sum of the three terms recovers the aggregated sector; the difference of
the classes $1$ and $2$ preserves the character's orientation; comparison
with the class divisible by three records the change of scale. These three
combinations are linearly independent and determine the three
coefficients $c_{r,n}$.

\section{Euler--Mascheroni and Stieltjes constants}

Let

\[
 C_r:=\operatorname*{FP}_{s=1}Z_r(s).
\]

The classical identities for the digamma function give
\cite{Apostol1976,Lagarias2013}

\begin{equation}
 C_0=\frac{\gamma-\log3}{3},
 \label{eq:C0-parte-finita}
\end{equation}

\begin{equation}
 C_1=\frac\gamma3+\frac{\log3}{6}+\frac{\pi}{6\sqrt3},
 \qquad
 C_2=\frac\gamma3+\frac{\log3}{6}-\frac{\pi}{6\sqrt3}.
 \label{eq:C1-C2-partes-finitas}
\end{equation}

\begin{corollary}
\label{cor:gamma-log3-L1}
The finite parts satisfy
\[
 \boxed{\gamma=C_0+C_1+C_2},
\]
\[
 \boxed{L(1,\chi_{-3})=C_1-C_2=\frac{\pi}{3\sqrt3}},
\]
\[
 \boxed{\log3=C_1+C_2-2C_0}.
\]
\end{corollary}

\begin{proof}
At order zero, the \cref{thm:descomposicion-laurent-triadica} gives
$a_0=c_{0,0}+c_{1,0}+c_{2,0}$. By definition,
$a_0=\gamma$ and $c_{r,0}=C_r$, proving the first equality. The
difference of the two oriented classes is
$C_1-C_2=L(1,\chi_{-3})$. Substituting
\eqref{eq:C0-parte-finita}--\eqref{eq:C1-C2-partes-finitas} yields the
other two identities.
\end{proof}

A single extraction of coefficients by Cauchy's integral
formula applied to \eqref{eq:coeficientes-Laurent-residuales} simultaneously determines
$\gamma_0,\ldots,\gamma_6$. The discretizations with 96 and
192 nodes agree up to order $10^{-87}$; the identity of the
coefficients, however, follows from the
\cref{thm:descomposicion-laurent-triadica} and not from that numerical comparison.

\section{Evaluation of Apéry's constant}

\begin{proposition}
\label{prop:evaluacion-apery}
The functional \eqref{eq:funcional-theta-mellin} satisfies
\[
 \boxed{
 \mathcal Z_3(3)=\zeta(3)
 =1.202056903159594285399738161511449\ldots}.
\]
\end{proposition}

\begin{proof}
The exact equality is the specialization $s=3$ of \eqref{eq:Z3-zeta}.
To certify any decimal prefix without modifying the functional,
consider $S_N=\sum_{n=1}^{N}n^{-3}$. The integral test gives
\[
 S_N+\frac1{2(N+1)^2}
 <\zeta(3)<
 S_N+\frac1{2N^2}.
\]
The resulting nested rational intervals determine the printed
digits.
\end{proof}

The irrationality of $\zeta(3)$ is Apéry's classical theorem
\cite{Apery1979}; it does not follow merely from evaluation of the functional.

\section{Glaisher--Kinkelin constant}

The spectral derivative at $s=-1$ defines

\begin{equation}
 \boxed{
 A_{\mathrm{GK}}
 =\exp\!\left(\frac1{12}-\mathcal Z_3'(-1)\right)}.
 \label{eq:glaisher-kinkelin-espectral}
\end{equation}

This is the formula for the regularized determinant associated with the
hyperfactorial product \cite{Vardi1988}. A five-point central differentiation,
evaluated at two scales, provides

\[
 A_{\mathrm{GK}}
 =1.2824271291006226368753425688697917\ldots.
\]

The term $1/12$ belongs to the regularization of the hyperfactorial and is not
introduced as a fitting parameter.

\section{Euler--Kronecker constant of the Eisenstein field}

The character $\chi_{-3}$ determines the imaginary quadratic field

\[
 F_{-3}=\QQ(\sqrt{-3}),
 \qquad
 \zeta_{F_{-3}}(s)=\zeta(s)L(s,\chi_{-3}).
\]

Its Euler--Kronecker constant is

\begin{equation}
 \boxed{
 \gamma_{F_{-3}}
 =\gamma+\frac{L'(1,\chi_{-3})}{L(1,\chi_{-3})}}.
 \label{eq:euler-kronecker-eisenstein}
\end{equation}

The residual coefficients of
\eqref{eq:coeficientes-Laurent-residuales} give

\[
 L(1,\chi_{-3})
 =0.6045997880780726168646927525473852\ldots,
\]
\[
 L'(1,\chi_{-3})
 =0.2226629869686015094866602627647443\ldots,
\]

and, by \eqref{eq:euler-kronecker-eisenstein},

\[
 \gamma_{F_{-3}}
 =0.9454972808716807032397499941581890\ldots.
\]

The choice of $F_{-3}$ comes from the residue character modulo three, not from
a search over the decimal values of constants of quadratic
fields.

\section{Local product associated with twin primes}

For a finite pattern $H=\{h_1,\ldots,h_k\}\subset\ZZ$, we define

\[
 \nu_p(H):=\left|\{-h_1,\ldots,-h_k\}\bmod p\right|,
 \qquad
 \sigma_p(H)=
 \frac{1-\nu_p(H)/p}{(1-1/p)^k}.
\]

For $H=\{0,2\}$, the local factors lead to

\begin{equation}
 C_2^{\mathrm{twin}}
 =\prod_{p>2}\frac{p(p-2)}{(p-1)^2}.
 \label{eq:constante-primos-gemelos}
\end{equation}

Let \(P_X\) be the product restricted to primes \(2<p\leq X\). For
\(p>X\), all factors belong to \((0,1)\); extending the tail
product from primes to all integers gives the telescoping
bound
\[
 \prod_{n>X}\left(1-\frac1{(n-1)^2}\right)
 =\frac{X-1}{X}
 \leq
 \prod_{p>X}\left(1-\frac1{(p-1)^2}\right)<1.
\]
The product \(P_{2\times10^6}\), evaluated with directed decimal rounding, and
this rational inequality for the tail give

\begin{equation}
 0.6601615071224244826403933551
 <C_2^{\mathrm{twin}}<
 0.6601618372035081248537566591.
 \label{eq:intervalo-primos-gemelos}
\end{equation}

The direct count up to $10^6$ contains 8169 pairs of twin primes. The
convergence of the local product and this finite census do not imply the
existence of infinitely many pairs; the latter is a distinct arithmetic
proposition.

\section{Meissel--Mertens constant and prime Hamiltonian}

The Meissel--Mertens constant is defined by

\begin{equation}
 B_1=\lim_{x\to\infty}
 \left(\sum_{p\leq x}\frac1p-\log\log x\right)
 \label{eq:meissel-mertens-def}
\end{equation}

and satisfies

\begin{equation}
 B_1=\gamma+
 \sum_p\left(\log\!\left(1-\frac1p\right)+\frac1p\right).
 \label{eq:meissel-mertens-euler}
\end{equation}

With the notation of \cref{pf84:C08}, we set
\[
 \mathcal F_P:=\Gamma_s(\mathfrak h_{\mathbb P}),
 \qquad
 H_P:=d\Gamma(H_{\mathbb P}).
\]
Neither the bosonic Fock space of prime modes nor its Hamiltonian is redefined here. The
operator already constructed has partition function

\begin{equation}
 \operatorname{Tr}_{\mathcal F_P}(e^{-sH_P})
 =\prod_p\frac1{1-p^{-s}}
 =\zeta(s),
 \qquad \Re s>1.
 \label{eq:hamiltoniano-primo-zeta}
\end{equation}

Let \(\mathcal H_P^{(1)}\) be the one-particle subspace, generated by
\(\{|p\rangle:p\text{ prime}\}\). Then
\[
 \operatorname{Tr}_{\mathcal H_P^{(1)}}\!\left(
 e^{-H_P}\mathbf 1_{(-\infty,\log x]}(H_P)\right)
 =\sum_{p\leq x}\frac1p.
\]
Thus, \(B_1\) is the finite part, as \(x\to\infty\), of this truncated spectral
trace after subtracting \(\log\log x\). The equality
\eqref{eq:hamiltoniano-primo-zeta} presupposes the Fock decomposition over
the primes; it therefore constitutes an operator representation of the Euler
product, not a deduction of primality from the triadic
Laplacian.

\section{Dependence on the operator}

The preceding evaluations share the operator
\eqref{eq:operador-espectral-Dm}, the positive trace
\eqref{eq:traza-positiva} and the transform
\eqref{eq:funcional-theta-mellin}. Other constants require distinct
dynamics. Catalan's constant is associated with the character modulo four
$\chi_{-4}$, whereas Khinchin's constant belongs to the Gauss
operator and its invariant measure. A prefix of any of them may
appear in a finite catalogue without the triadic functional selecting it.
This operator dependence distinguishes spectral evaluation from the
purely combinatorial capacity for encoding.
